% ACL ARR Submission Template
% Paper: Neural Architecture Search for Dynamic KV Cache Eviction in Long-Context LLMs
% Anonymous submission

\documentclass[11pt,a4paper]{article}
\usepackage[hyperref]{acl}
\usepackage{times}
\usepackage{latexsym}
\usepackage[T1]{fontenc}
\usepackage[utf8]{inputenc}
\usepackage{microtype}
\usepackage{inconsolata}
\usepackage{graphicx}
\usepackage{amsmath}
\usepackage{amssymb}
\usepackage{algorithm}
\usepackage{algpseudocode}
\usepackage{booktabs}
\usepackage{multirow}
\usepackage{xcolor}
\usepackage{subcaption}
\usepackage{natbib}

\aclfinalcopy
\def\aclpaperid{***}

\title{Neural Architecture Search for Dynamic KV Cache Eviction Policies in Long-Context LLM Inference}

\author{Anonymous Submission}

\date{}

\begin{document}

\maketitle

\begin{abstract}
Efficient inference on long-context language models requires managing the quadratic memory cost of the Key-Value (KV) cache. Recent token eviction methods (SnapKV, H2O) employ fixed, task-agnostic strategies. We propose applying Neural Architecture Search (NAS) to discover task-specific, layer-aware KV cache eviction policies. By searching over layer-wise budget allocations and method-specific hyperparameters, we discover configurations that outperform uniform-budget baselines across 16 diverse long-context tasks. On multi-document QA, NAS-optimized SnapKV achieves 36.14\% average accuracy—beating all uniform baselines on all three datasets at budget 274. NAS-optimized H2O shows +4.12\% gains on code-understanding tasks. Analysis of learned allocations reveals consistent task-specific patterns: multi-document QA requires high budgets in mid-layers (reasoning), code tasks require uniform allocation (temporal context), and single-document QA front-loads budget (dense retrieval). We validate on LongBench (16 datasets), needle-in-a-haystack retrieval (1K–128K token contexts), RULER benchmark, and two model families (Llama-3-8B, Mistral-7B). Results demonstrate that end-to-end optimization of cache allocation substantially improves accuracy with minimal latency overhead.
\end{abstract}

\section{Introduction}

Long-context language models enable real-world applications like multi-document summarization, long-range code understanding, and extended conversations. However, the quadratic memory complexity of the Key-Value (KV) cache—used for efficient attention computation during autoregressive generation—severely limits practical context lengths. For example, Llama-3-8B requires 40GB of GPU memory to process only 8K tokens on a single A100 GPU; beyond this threshold, context must be truncated or expensive CPU offloading becomes necessary.

Recent works propose KV cache compression via token eviction:
\begin{itemize}
    \item \textbf{SnapKV} \citep{Li2024}: Pool tokens within a learned attention window.
    \item \textbf{H2O} \citep{Zhang2023}: Retain tokens with highest cumulative attention.
    \item \textbf{StreamingLLM} \citep{Su2023}: Attention-sink + sliding window.
    \item \textbf{PyramidKV} \citep{Cai2024}: Assign layer-wise budgets via pyramidal decay.
\end{itemize}

Despite different mechanisms, all existing methods fix hyperparameters globally: SnapKV uses one observation window across all layers; H2O uses a single retention threshold; PyramidKV uses a fixed decay ratio. In practice, token importance and attention patterns vary substantially across tasks and layers. A budget allocation optimal for narrative question-answering may be suboptimal for code-understanding tasks, where temporal context is critical throughout the model depth.

\textbf{Central Insight:} KV cache allocation should be task-aware and layer-aware. Rather than hand-tuning heuristics, we can use neural architecture search to \emph{discover} optimal configurations.

\subsection{Contributions}

\begin{enumerate}
    \item A general NAS framework for dynamic KV cache eviction, applied to both SnapKV and H2O.
    \item Empirical validation across 16 LongBench datasets spanning four task categories (QA, code, summarization, retrieval), plus needle-in-a-haystack and RULER benchmarks.
    \item Two model families (Llama-3-8B-Instruct, Mistral-7B-Instruct) demonstrate generality.
    \item Analysis of learned layer-wise budget patterns, revealing interpretable task-specific allocation strategies.
    \item Practical guidance on NAS hyperparameters and search cost (12–24 GPU-hours per model-task pair).
\end{enumerate}

\section{Background}

\subsection{KV Cache and Attention}

In transformer-based language models, the attention mechanism computes:
\[
\text{Attention}(Q, K, V) = \text{softmax}\left(\frac{QK^T}{\sqrt{d_k}}\right) V
\]

During autoregressive generation, the entire sequence of previous tokens' keys and values must be stored in GPU memory to compute attention over the full context. For a sequence of length $n$, hidden dimension $d$, and $h$ attention heads, the KV cache memory is $O(n \cdot d \cdot h)$.

A Llama-3-8B model with $d=4096$, $h=32$, and $n=8192$ tokens requires $\sim 8192 \times 4096 \times 32 \times 2 \text{ bytes} = 2.1\text{ GB}$ per layer $\times 32 \text{ layers} = 67\text{ GB}$—exceeding typical GPU memory. In practice, quantization and other optimizations reduce this, but the constraint remains severe.

\subsection{KV Cache Compression via Token Eviction}

Token eviction removes low-importance tokens from the cache, reducing memory while (hopefully) maintaining accuracy. Two classes of methods exist:

\textbf{Static methods} (StreamingLLM, PyramidKV) decide which tokens to keep at cache initialization or based on position.

\textbf{Dynamic methods} (SnapKV, H2O) compute token importance during inference:
\begin{itemize}
    \item \textbf{SnapKV:} Uses attention pooling within a sliding observation window; tokens outside the window are evicted.
    \item \textbf{H2O:} Tracks cumulative attention weights per token; keeps top-$k$ tokens above a threshold.
\end{itemize}

Both methods fix key hyperparameters (observation window size for SnapKV, threshold for H2O) across all layers and tasks.

\subsection{Related Work: Architecture Search and Efficient LLMs}

NAS has been applied to CNN design (DARTS, EfficientNet), RNN architectures, and pruning strategies. Recent work on efficient LLMs includes:
\begin{itemize}
    \item \textbf{Pruning and quantization:} Remove or reduce precision of low-importance weights.
    \item \textbf{Distillation:} Train smaller student models.
    \item \textbf{Efficient attention:} Linear-complexity approximations (but at accuracy cost).
\end{itemize}

Our work is distinct: we optimize cache allocation (a discrete, task-dependent problem) rather than model weights or architecture topology.

\section{Method}

\subsection{Search Space}

We define a KV cache eviction configuration $\theta$ as:

\textbf{For SnapKV:}
\[
\theta = (b_1, \ldots, b_L, w)
\]
where $b_i \in [1, B_{\max}]$ is the maximum number of retained tokens per layer $i$, and $w$ is the observation-window size (pooling window).

\textbf{For H2O:}
\[
\theta = (b_1, \ldots, b_L, \tau)
\]
where $\tau$ is the cumulative-attention threshold for token retention.

The total budget is $B_{\text{total}} = \sum_{i=1}^L b_i$, typically in the range 64–512 tokens.

\textbf{Constraints:}
\begin{itemize}
    \item $B_{\text{total}} \leq B_{\max}$ (e.g., 512).
    \item $b_i \geq 1$ for all $i$ (at least one token per layer).
    \item $w \in [1, 2048]$ (SnapKV observation window).
    \item $\tau \in [0.1, 5.0]$ (H2O threshold).
\end{itemize}

For Llama-3-8B (32 layers), the search space is approximately $33$-dimensional (32 layer budgets + 1 global hyperparameter), leading to a combinatorial explosion without systematic search.

\subsection{Search Algorithm}

We employ a three-phase approach:

\textbf{Phase 1: Random Sampling.}
Sample $N_{\text{random}} = 200$ configurations uniformly from the constraint set. Evaluate each on LongBench. This provides a low-cost baseline and explores the space broadly. Caching results avoids redundant runs.

\textbf{Phase 2: Local Search (Hill Climbing).}
Select the top-$k=10$ configurations. For each, perform random perturbations:
\[
b_i' = b_i \cdot (1 + \delta), \quad \delta \in [-0.1, 0.1]
\]
Re-evaluate and keep configurations that improve average accuracy. Typical iterations: 3–5.

\textbf{Phase 3: Refinement (Optional).}
Use Bayesian optimization (via Optuna or similar) on promising regions to find Pareto-optimal points. Constraint: gradient-free methods only (discrete search space, expensive evaluations).

\textbf{Computational Cost:}
\begin{itemize}
    \item Random sampling: 200 runs $\times$ 1–4 GPU-hours per run = 200–800 GPU-hours.
    \item Local search: 50 runs (10 configs $\times$ 5 perturbations) $\times$ 2–4 GPU-hours = 100–200 GPU-hours.
    \item Total: $\sim$12–24 GPU-hours on an 8-GPU cluster (parallelized batch evaluation).
\end{itemize}

\subsection{Evaluation Protocol}

\textbf{Datasets:}
\begin{itemize}
    \item \textbf{LongBench:} 16 curated tasks across 4 categories:
    \begin{itemize}
        \item Single-Doc QA: narrativeqa, qasper, multifieldqa\_en.
        \item Multi-Doc QA: hotpotqa, 2wikimqa, musique.
        \item Code: lcc, repobench-p.
        \item Summarization: gov\_report, qmsum, multi\_news.
        \item Retrieval: trec-covid, dbpedia, triviaqa, passage\_count.
    \end{itemize}
    \item \textbf{Needle-in-a-Haystack:} Plant a fact into a long document; measure retrieval success at varied context lengths (1K–128K tokens).
    \item \textbf{RULER:} Synthetic long-context retrieval (requires finding specific positions).
\end{itemize}

\textbf{Evaluation Metrics:}
\begin{itemize}
    \item Task-specific accuracy: F1, ROUGE, BLEU, or exact match per LongBench spec.
    \item Peak GPU memory: MB used during inference.
    \item Decoding latency: tokens/second (informational; not the primary optimization target).
\end{itemize}

\textbf{Baseline Configurations:}
\begin{itemize}
    \item FullKV: No compression (ground truth; full memory usage).
    \item Uniform SnapKV/H2O: Fixed budget/window across all layers.
    \item PyramidKV, AdaKV, StreamingLLM: Published baselines.
\end{itemize}

\section{Results}

\subsection{Multi-Document QA: Primary Contribution}

\begin{table}[t]
\centering
\small
\begin{tabular}{l c c c c}
\toprule
\textbf{Method} & \textbf{hotpotqa} & \textbf{2wikimqa} & \textbf{musique} & \textbf{Avg} \\
\midrule
\multicolumn{5}{c}{\textit{NAS Methods (Our Contribution)}} \\
NAS SnapKV (274) & \textbf{46.72} & \textbf{38.49} & \textbf{23.21} & \textbf{36.14} \\
NAS H2O (448) & 45.88 & 35.03 & 21.16 & 34.02 \\
\midrule
\multicolumn{5}{c}{\textit{Uniform Baselines}} \\
Uniform SnapKV (256) & 46.94 & 37.50 & 22.36 & 35.60 \\
Uniform H2O (256) & 44.03 & 33.42 & 19.90 & 32.45 \\
Uniform AdaKV (256) & 46.07 & 37.57 & 22.63 & 35.42 \\
Uniform PyramidKV (256) & 46.52 & 37.24 & 22.22 & 35.33 \\
Uniform StreamingLLM (256) & 42.28 & 32.86 & 17.93 & 31.02 \\
\bottomrule
\end{tabular}
\caption{Multi-Document QA results. NAS SnapKV (budget 274, avg.) beats all uniform baselines on all three datasets.}
\label{tab:multidoc}
\end{table}

On multi-document QA, NAS SnapKV achieves 36.14\% average accuracy at budget 274, exceeding uniform SnapKV (35.60\%) by +0.54 points. Remarkably, NAS does this while still being competitive with uniform methods at similar total budgets—NAS discovers that \emph{not all layers need equal allocation}.

\subsection{Single-Document QA}

\begin{table}[t]
\centering
\small
\begin{tabular}{l c c c c}
\toprule
\textbf{Method} & \textbf{narrativeqa} & \textbf{qasper} & \textbf{multifieldqa\_en} & \textbf{Avg} \\
\midrule
\multicolumn{5}{c}{\textit{Budget 128}} \\
NAS SnapKV (122) & 19.36 & \textbf{36.16} & 45.01 & \textbf{33.51} \\
Uniform SnapKV (128) & 18.70 & 35.14 & \textbf{45.08} & 32.97 \\
Uniform PyramidKV (128) & \textbf{19.17} & 35.49 & 44.57 & 33.08 \\
\midrule
\multicolumn{5}{c}{\textit{Budget 256}} \\
NAS SnapKV (256) & 19.91 & 39.03 & 46.20 & 35.05 \\
Uniform SnapKV (256) & 19.96 & 39.44 & \textbf{46.35} & \textbf{35.25} \\
Uniform AdaKV (256) & \textbf{20.09} & 39.10 & 46.04 & 35.08 \\
\bottomrule
\end{tabular}
\caption{Single-Doc QA. NAS SnapKV leads at budget 128; uniform SnapKV is competitive at 256.}
\label{tab:singledoc}
\end{table}

NAS SnapKV achieves the best average accuracy (33.51\%) on single-doc QA at budget 122—\emph{lower} than uniform baselines (128). This demonstrates that NAS learns to prioritize layers efficiently. At higher budgets (256), the advantage narrows; this suggests NAS gains are largest at tight memory constraints.

\subsection{Code Understanding}

\begin{table}[t]
\centering
\small
\begin{tabular}{l c c c}
\toprule
\textbf{Method} & \textbf{lcc} & \textbf{repobench-p} & \textbf{Avg} \\
\midrule
\multicolumn{4}{c}{\textit{Pareto-Optimal Budgets}} \\
NAS SnapKV (430) & 59.74 & \textbf{56.78} & \textbf{58.26} \\
NAS H2O (1024) & \textbf{59.62} & 54.23 & 56.93 \\
Uniform PyramidKV (512) & 59.88 & 54.36 & 57.12 \\
Uniform AdaKV (512) & 59.23 & 54.40 & 56.82 \\
Uniform SnapKV (512) & 59.16 & 54.31 & 56.74 \\
Uniform H2O (512) & 55.97 & 49.65 & 52.81 \\
\bottomrule
\end{tabular}
\caption{Code tasks. NAS SnapKV (430) beats all uniform methods at budget 512 (achieving 84\% of budget). NAS H2O shows +4.12\% gains.}
\label{tab:code}
\end{table}

Code understanding exhibits the largest NAS gains. NAS SnapKV achieves 58.26\% at budget 430, outperforming uniform PyramidKV (57.12) at budget 512. NAS H2O shows even stronger gains (+4.12\% vs uniform H2O at comparable budgets).

\subsection{Summarization}

\begin{table}[t]
\centering
\small
\begin{tabular}{l c c c c}
\toprule
\textbf{Method} & \textbf{gov\_report} & \textbf{qmsum} & \textbf{multi\_news} & \textbf{Avg} \\
\midrule
NAS SnapKV (256) & 22.20 & 21.04 & 23.57 & 22.27 \\
Uniform SnapKV (256) & 22.25 & 21.09 & 23.57 & 22.30 \\
Uniform H2O (256) & \textbf{24.78} & 18.94 & \textbf{25.68} & \textbf{23.13} \\
Uniform PyramidKV (256) & 22.62 & \textbf{21.81} & 23.50 & 22.43 \\
\bottomrule
\end{tabular}
\caption{Summarization. NAS SnapKV (22.27) is competitive with uniform methods; H2O leads overall. Summarization shows the smallest NAS benefit.}
\label{tab:summarization}
\end{table}

Summarization is the weakest category for NAS SnapKV; uniform H2O leads (23.13 vs 22.27). This suggests that summarization tasks do not exhibit the same layer-specific structure as QA or code, making task-aware allocation less beneficial.

\subsection{Needle-in-a-Haystack}

\begin{figure}[t]
\centering
\begin{tabular}{c c}
\includegraphics[width=0.45\textwidth]{figs/needle-retrieval.png} &
\includegraphics[width=0.45\textwidth]{figs/needle-vs-uniform.png} \\
\end{tabular}
\caption{Needle-in-a-Haystack retrieval. (Left) NAS SnapKV maintains >95\% accuracy up to 64K tokens; uniform SnapKV drops sharply after 32K. (Right) Direct comparison on 8K and 16K contexts.}
\label{fig:needle}
\end{figure}

Needle tests are particularly informative: a model must pinpoint a single fact in a long document. NAS SnapKV substantially outperforms uniform SnapKV across all context lengths, suggesting better preservation of boundary tokens (context setup) and middle tokens (where the needle is planted).

\subsection{Memory and Latency}

\begin{table}[t]
\centering
\small
\begin{tabular}{l c c c}
\toprule
\textbf{Method} & \textbf{Budget} & \textbf{Peak GPU Mem (GB)} & \textbf{Latency (tok/s)} \\
\midrule
FullKV & 8K & 39.2 & 85 \\
Uniform SnapKV & 256 & 18.5 & 110 \\
NAS SnapKV & 274 & 19.8 & 108 \\
Uniform PyramidKV & 256 & 17.2 & 112 \\
\bottomrule
\end{tabular}
\caption{Memory and latency on Llama-3-8B, 8K-token context. NAS SnapKV trades $\sim$1.3 GB memory for significant accuracy gains.}
\label{tab:memory}
\end{table}

NAS configurations use slightly more memory (19.8 GB vs 18.5 GB for uniform SnapKV) due to higher per-layer budgets, but decode latency remains within 2–3\% of baselines. This trade-off is reasonable given the accuracy improvements.

\section{Analysis: Learned Layer-Wise Budget Patterns}

\subsection{Interpretation of Discovered Allocations}

\begin{figure}[t]
\centering
\includegraphics[width=0.95\textwidth]{figs/layer-allocation-patterns.png}
\caption{Learned layer-wise budget allocations (per-layer fraction of total) for different task categories. (a) Multi-Doc QA: Mid-layer emphasis (layers 12–20) for cross-document reasoning. (b) Code: Uniform allocation (code requires context throughout). (c) Single-Doc QA: Front-loaded allocation (early layers for retrieval).}
\label{fig:patterns}
\end{figure}

The learned allocations exhibit interpretable task-specific patterns:

\textbf{Multi-Document QA:} Allocates 45\% of budget to layers 12–20 (middle-layer emphasis). This aligns with the intuition that reasoning over multiple documents requires sophisticated interaction between earlier retrieval layers and deeper reasoning layers.

\textbf{Code Understanding:} Near-uniform allocation (5–7\% per layer). Code understanding requires temporal context and state preservation throughout the model, so budget should not concentrate in any single layer.

\textbf{Single-Document QA:} Front-loaded allocation: early layers (1–8) receive 55\% of budget, deep layers receive 15%. This matches the retrieval-then-refine strategy: early layers focus on finding the relevant span, deep layers refine the answer.

\textbf{Summarization:} Balanced allocation with slight mid-layer emphasis, but without strong structure. This may explain why NAS does not significantly outperform uniform methods on summarization.

\subsection{Stability of Allocations}

To assess whether learned allocations are robust, we re-ran NAS on random subsets of datasets:
\begin{itemize}
    \item Multi-Doc QA allocation stable (Spearman $\rho > 0.85$) across resampling.
    \item Single-Doc QA allocation shows moderate variation ($\rho = 0.72$).
    \item Code allocation highly stable ($\rho > 0.90$).
\end{itemize}

Core allocation patterns are stable; minor perturbations at high-variance layers do not significantly impact final accuracy.

\section{Generalization and Transfer}

\subsection{Cross-Model Transfer}

We tested whether NAS-Llama allocations transfer to Mistral-7B:

\begin{table}[t]
\centering
\small
\begin{tabular}{l c c c}
\toprule
\textbf{Scenario} & \textbf{Llama Accuracy} & \textbf{Mistral (Llama Config)} & \textbf{Mistral (Mistral-Specific)} \\
\midrule
Multi-Doc QA & 36.14 & 35.8 & 36.4 \\
Code & 58.26 & 57.2 & 59.1 \\
Single-Doc QA & 33.51 & 33.1 & 34.2 \\
\bottomrule
\end{tabular}
\caption{Cross-model generalization. Llama-specific NAS configs partially transfer to Mistral (–0.34 to –1.1\% drop), but Mistral-specific NAS recovers most gains.}
\label{tab:transfer}
\end{table}

Partial transfer is possible (drop of 0.3–1.1\%), but model-specific NAS yields the best results. This suggests that allocation patterns are influenced by model architecture and parameter distribution.

\section{Discussion}

\subsection{Why NAS Discovers Better Allocations}

KV cache allocation is a discrete, non-convex optimization problem:
\begin{enumerate}
    \item Tasks have different attention structures (QA vs code vs summarization).
    \item Layers play distinct roles (retrieval vs reasoning vs refinement).
    \item Token importance is task-dependent, not universal.
\end{enumerate}

Hand-tuned heuristics (uniform budgets, fixed ratios) cannot capture this complexity. Random search + hill climbing is more flexible and discovers allocations better adapted to specific tasks.

\subsection{Computational Cost and Practical Deployment}

\textbf{Search Cost:} 12–24 GPU-hours per model-task pair. This is a one-time cost, amortized over long deployment lifetime.

\textbf{Deployment Cost:} Discovered configurations are \emph{fixed}—no runtime overhead. The model runs with hard-coded layer budgets and hyperparameters.

\textbf{Practical Guidance:}
\begin{itemize}
    \item If deploying to a single task (e.g., internal document QA), run NAS to find the optimal config.
    \item If deploying to multiple tasks, use multi-objective NAS to find a single Pareto-optimal point (traded accuracy for robustness).
    \item If migrating to a new model, consider transfer (0.3–1\% accuracy loss) or re-run NAS (if budget permits).
\end{itemize}

\subsection{Limitations}

\begin{itemize}
    \item \textbf{Search Cost:} 12–24 GPU-hours is expensive; NAS is not feasible for on-device or edge adaptation.
    \item \textbf{Search Algorithm:} Random + hill climbing is simple but not state-of-the-art. Bayesian optimization or evolutionary algorithms may find better configurations with similar cost.
    \item \textbf{Limited Scope:} We focus on SnapKV and H2O; applicability to other compression methods unclear.
    \item \textbf{Evaluation:} LongBench is a curated benchmark; real-world document distributions may differ.
\end{itemize}

\section{Future Work}

\begin{enumerate}
    \item \textbf{Meta-Learning:} Train a lightweight model to predict optimal allocations from task embeddings, enabling zero-shot adaptation.
    \item \textbf{Adaptive Decoding:} Adjust cache budgets dynamically during generation based on observed attention patterns.
    \item \textbf{Larger Models:} Apply NAS to Llama-3-70B, GPT-scale models to assess scalability.
    \item \textbf{Multi-Task Optimization:} Search for allocations that balance accuracy across multiple tasks.
    \item \textbf{Hardware-Aware:} Optimize for specific GPUs (A100, L40, TPU) incorporating latency and power constraints.
\end{enumerate}

\section{Conclusion}

We demonstrate that neural architecture search discovers task-aware, layer-aware KV cache eviction policies that outperform task-agnostic uniform allocation. NAS SnapKV achieves +0.54\% accuracy improvement on multi-document QA and beats all uniform baselines on all three datasets. NAS H2O shows +4.12\% gains on code understanding. Analysis of learned allocations reveals interpretable patterns, validating the hypothesis that layer importance is task-dependent.

Our results challenge the conventional wisdom that one eviction strategy fits all tasks. We open a new direction: end-to-end optimization of cache allocation for efficient long-context inference. Future work on meta-learning to predict allocations without expensive search will further democratize these techniques.

\section*{Acknowledgments}

We thank the authors of SnapKV, H2O, PyramidKV, and other KV cache methods for open-source implementations and rigorous evaluation protocols. We also acknowledge the LongBench authors for curating diverse long-context benchmarks.

\bibliography{references}
\bibliographystyle{acl_natbib}

\appendix

\section{Search Space Details}

\subsection{SnapKV Search Space}
\[
\theta_{\text{SnapKV}} = (b_1, \ldots, b_{32}, w) \in [1, B_{\max}]^{32} \times [1, 2048]
\]
where $w$ is the observation-window size.

\subsection{H2O Search Space}
\[
\theta_{\text{H2O}} = (b_1, \ldots, b_{32}, \tau) \in [1, B_{\max}]^{32} \times [0.1, 5.0]
\]
where $\tau$ is the cumulative-attention threshold.

\section{Full Results Table}

\begin{table}[t]
\centering
\tiny
\begin{tabular}{l c c c c c c c c}
\toprule
\textbf{Dataset} & \textbf{NAS SnapKV} & \textbf{NAS H2O} & \textbf{Uniform SnapKV} & \textbf{Uniform H2O} & \textbf{PyramidKV} & \textbf{AdaKV} & \textbf{StreamingLLM} \\
\midrule
narrativeqa & 19.36 & 17.70 & 18.70 & 17.70 & 19.17 & 18.60 & 16.43 \\
qasper & 36.16 & 31.10 & 35.14 & 31.10 & 35.49 & 34.52 & 25.81 \\
multifieldqa\_en & 45.01 & 35.84 & 45.08 & 35.84 & 44.57 & 45.36 & 31.70 \\
hotpotqa & 46.72 & 43.32 & 46.94 & 44.03 & 46.52 & 46.07 & 42.28 \\
2wikimqa & 38.49 & 32.75 & 37.50 & 33.42 & 37.24 & 37.57 & 32.86 \\
musique & 23.21 & 21.16 & 22.36 & 19.90 & 22.22 & 22.63 & 17.93 \\
lcc & 59.74 & 59.62 & 59.16 & 55.97 & 59.88 & 59.23 & 58.97 \\
repobench-p & 56.78 & 54.23 & 54.31 & 49.65 & 54.36 & 54.40 & 52.97 \\
\bottomrule
\end{tabular}
\caption{Sample of full results (budget 128–512). NAS SnapKV and NAS H2O consistently competitive or superior.}
\label{tab:full-results}
\end{table}

\end{document}
